% ============================================================================
%  Chapter 17 -- Verm1, Part 3: A Body, Rules and States         (Lecture 17)
% ============================================================================
\chapter{Verm1, Part 3: A Body, Rules and States}
\label{ch:verm1p3}

\epigraphline{Show me your tables, and I won't usually need your flowcharts;
they'll be obvious.}{Fred Brooks, \emph{The Mythical Man-Month}}

\begin{chapgoals}
By the end of this chapter you will have
\begin{itemize}
\item a worm with a body that follows the head and grows by one cell at every meal;
\item a game that can be lost: the worm must not bite itself;
\item three screens, a title, the game and a game over, linked by button \key{A};
\item food that never appears on the worm;
\item a best score, kept as long as the cartridge stays in memory.
\end{itemize}
\end{chapgoals}

In this third tutorial memory enters the game: a ring of bytes, a map of the
world, a machine with three states. With ``data structures'' in the place of
``tables'', the sentence of Brooks quoted above is the motto of the chapter:
once the data are designed well, the code is a transcription.
Table~\ref{tab:v3:from} lists what is taken from Chapters \ref{ch:sequential}
to~\ref{ch:storage}. The frame is the clock of the machine: \code{update}
computes the next state from the present state and from the inputs. The
reference file is \fname{verm1\_part3.S} (Appendix~\ref{app:verm1}).

\begin{table}[htb]
\centering\small
\caption{What Part 3 takes from the chapters on state and memory.}\label{tab:v3:from}
\begin{tabularx}{\textwidth}{@{}lLL@{}}
\toprule
\textbf{From} & \textbf{Idea} & \textbf{Where it appears}\\
\midrule
Chapter \ref{ch:sequential} & a finite-state machine: states, inputs, transitions & \texttt{TITLE}, \texttt{PLAY}, \texttt{OVER}\\
Chapter \ref{ch:sequential} & an edge is a change, not a level & \key{A} pressed now and not one frame ago\\
Chapter \ref{ch:registers} & a register keeps a value for the next clock & \code{prev\_btn}: the buttons of the last frame\\
Chapter \ref{ch:memory} & memory is an array of bytes: address = base + index & the body, the map of the world\\
Chapter \ref{ch:storage} & what stays in the file and what lives in RAM & \code{.rodata} and \code{.data}; the best score\\
\bottomrule
\end{tabularx}
\end{table}

\begin{figure}[htb]
\centering
\begin{subfigure}[b]{0.32\textwidth}\includegraphics[width=\textwidth]{verm1_p3_title.png}\caption{\texttt{TITLE}}\end{subfigure}\hfill
\begin{subfigure}[b]{0.32\textwidth}\includegraphics[width=\textwidth]{verm1_p3_play.png}\caption{\texttt{PLAY}: seven cells after four meals}\label{fig:v3:play}\end{subfigure}\hfill
\begin{subfigure}[b]{0.32\textwidth}\includegraphics[width=\textwidth]{verm1_p3_over.png}\caption{\texttt{OVER}}\label{fig:v3:over}\end{subfigure}
\caption{Where this chapter is going: three screens, one program.}
\label{fig:v3:screens}
\end{figure}

Start from a Part~2 that runs; if yours does not, copy
\fname{verm1\_part2.S}. The working loop is that of
Listing~\ref{lst:v1:build}. In this part the file changes a great deal: build
after every \emph{Run it}, not only at the end.

% ----------------------------------------------------------------------------
\section{The design, before the code}

Two data structures and one diagram are designed first.

\subsection{The body is a ring of 256 entries}

The body is kept in two arrays of bytes, \code{body\_x} and \code{body\_y}:
entry $k$ holds the cell of one segment. Nothing is ever shifted: a step
writes one new entry and moves an index (Figure~\ref{fig:v3:ring}). Two
numbers describe the worm: \code{head\_i}, the entry of the head, and
\code{length}, the number of entries in use. The tail is the entry
\[ \code{head\_i}-\code{length}+1 . \]
After entry 255 comes entry 0, and the wrap is one instruction: the index AND
255. The structure is a \term{ring}, or circular buffer: the structure
behind every first-in, first-out queue.

\begin{figure}[htb]
\centering
\begin{tikzpicture}[x=7mm,y=7mm]
  \foreach \row/\h/\len/\cap in {0/2/3/{A new game: three cells, the head in entry 2},
                                 1/3/3/{One step: write the new head in entry 3; the tail is now entry 1},
                                 2/4/4/{One step with a meal: the head advances, the tail stays}}{
    \begin{scope}[yshift=-\row*19mm]
      \foreach \k in {0,...,15}{
        \pgfmathtruncatemacro{\tail}{\h-\len+1}
        \ifnum\k=\h \def\col{green!75!black}\else\ifnum\k<\tail \def\col{prgpanel}\else\ifnum\k>\h \def\col{prgpanel}\else\def\col{prggreen!55}\fi\fi\fi
        \draw[prgsteel,fill=\col] (\k,0) rectangle ++(1,1);
        \node[font=\ttfamily\scriptsize,text=prgcomment] at (\k+0.5,1.3) {\k};
      }
      \node[font=\ttfamily\scriptsize] at (16.6,0.5) {\ldots 255};
      \node[lbl,anchor=west] at (0,-0.45) {\cap};
      \draw[->,thick,prgred] (\h+0.5,2.1) node[above,font=\ttfamily\scriptsize,text=prgred]{head\_i} -- (\h+0.5,1.55);
    \end{scope}
  }
\end{tikzpicture}
\caption{The first sixteen entries of the ring in three successive moments.}
\label{fig:v3:ring}
\end{figure}

Moving a worm of 200 cells therefore costs the same as moving one of 3. The
size 256 is chosen because it is a power of two, so that the wrap is a mask.

\subsection{The world is a map of 960 bytes}

Does the new head touch the body? Looking at every segment would cost up to
255 comparisons. Instead the program keeps a \emph{map}: one byte per cell, 1
if a segment is there and 0 if not. The cell in column $c$ of row $r$ is the
byte number
\[ r\times 40+c,\qquad\text{from 0 to 959}, \]
and its address is $\code{grid}+r\times 40+c$: a two-dimensional world in a
one-dimensional memory, the address computation of Chapter~\ref{ch:memory}. A
step sets one byte, that of the new head, and clears one, that of the old
tail. The collision test is a single load, whatever the length of the worm.

It is the classic exchange of memory for time. The map and the ring say the
same thing in two ways, and the program must keep them consistent: every write
to one is paired with a write to the other.

\subsection{The game is a machine with three states}

\begin{figure}[htb]
\centering
\begin{tikzpicture}[->,shorten >=1pt,every edge/.append style={font=\sffamily\footnotesize}]
  \node[st,fill=prgsteel,text=white,minimum size=15mm] (t) at (0,0) {TITLE};
  \node[st,fill=prgaccent,text=white,minimum size=15mm] (p) at (4.6,0) {PLAY};
  \node[st,fill=prgred,text=white,minimum size=15mm] (o) at (9.2,0) {OVER};
  \path (t) edge node[above]{\textbf{A} pressed} node[below]{\scriptsize\texttt{new\_game}} (p)
        (p) edge node[above]{head on the body} node[below]{\scriptsize update \texttt{best}} (o)
        (o) edge[bend left=32] node[below]{\textbf{A} pressed} (t);
\end{tikzpicture}
\caption{The state diagram of \vermone. In every other case the state does not change.}
\label{fig:v3:fsm}
\end{figure}

The state is one word in memory (Figure~\ref{fig:v3:fsm}). \code{update}
looks at the state first and then at the inputs. And ``\key{A} pressed''
means: down now, and up one frame ago. Holding \key{A} must not skip the
title: what is wanted is an edge, not a level.

% ----------------------------------------------------------------------------
\section{Step 1: the memory of the game}

Replace the constants, the \code{.rodata} and the \code{.data} of Part~2 with
Listings \ref{lst:v3:const}, \ref{lst:v3:ro} and~\ref{lst:v3:data}.

\asmfile[caption={Step 1 \textbf{(changed)}: constants.},label=lst:v3:const]{code/l17/s11.S}

There is a darker green, \hex{03e0}, for the body; \code{BTN\_A} is bit~4 of
the word of buttons; the states are numbers, 0, 1 and 2.

\asmfile[caption={Step 1 \textbf{(changed)}: \code{.rodata}, what stays in the file.},label=lst:v3:ro]{code/l17/s12.S}

Three strings and the two tables of Part~2: read-only data, to which the
program never writes, and which on the board could stay in flash.

\asmfile[caption={Step 1 \textbf{(changed)}: \code{.data}, what lives in RAM.},label=lst:v3:data]{code/l17/s13.S}

\code{head\_x} and \code{head\_y} are gone: the head is in the ring. The
directive \code{.space} reserves bytes without naming each one:
$256+256+960$ bytes, the worm and its world. Everything here is lost when the
cartridge is unloaded (Chapter~\ref{ch:storage}): \code{best} survives a game,
not a power-off. A best score that survives would need permanent memory,
which a portable cartridge does not have: an honest limit.

% ----------------------------------------------------------------------------
\section{Step 2: three states}

The machine is built first, with functions that do nothing yet.

\subsection{\texttt{init}: start on the title}

Replace \code{verm1\_init} with Listing~\ref{lst:v3:init}, and delete from
the file \code{verm1\_update}, \code{step} and \code{verm1\_draw} of Part~2:
the next listings replace them. Keep \code{place\_food}, \code{eat\_sound}
and \code{draw\_cell}.

\asmfile[caption={Step 2 \textbf{(changed)}: \code{verm1\_init}.},label=lst:v3:init]{code/l17/s15.S}

\code{init} only chooses the first state; the game itself is prepared by
\code{new\_game}.

\subsection{\texttt{update}: buttons, and their edges}

\asmfile[caption={Step 2 \textbf{(changed)}: \code{verm1\_update}, first half.},label=lst:v3:edge]{code/l17/s16.S}

In Listing~\ref{lst:v3:edge} \reg{s0} holds the buttons now and
\code{prev\_btn} the buttons one frame ago. The variable \code{prev\_btn} is
a register in the sense of Chapter~\ref{ch:registers}: it delays its input by
one clock, and here the clock is the frame. The expression ``now AND NOT
before'' is 1 only where a press has begun:
\[ \reg{s1}=\reg{s0}\cdot\nott{\code{prev\_btn}} . \]
It is the edge detector of Chapter~\ref{ch:sequential}: a flip-flop and an AND
gate with one inverted input. \reg{s0} and \reg{s1} must survive the calls,
and are saved in the frame.

\subsection{\texttt{update}: one branch per state}

\asmfile[caption={Step 2 \textbf{(changed)}: \code{verm1\_update}, second half.},label=lst:v3:states]{code/l17/s17.S}

Listing~\ref{lst:v3:states} is the transition table, line by line. The state
is read first: it decides who looks at the inputs. \texttt{TITLE} and
\texttt{OVER} react only to a new press of \key{A}; \texttt{PLAY} passes the
buttons that are held to \code{play\_update}. The transition from
\texttt{OVER} to \texttt{TITLE} only stores the new state; that from
\texttt{TITLE} to \texttt{PLAY} calls \code{new\_game}, which stores
\texttt{PLAY} itself.

\subsection{\texttt{draw}: what to show in each state}

\asmfile[caption={Step 2 \textbf{(changed)}: \code{verm1\_draw}, first half.},label=lst:v3:draw1]{code/l17/s18.S}

The title line is always drawn (Listing~\ref{lst:v3:draw1}). In
\texttt{TITLE} only the invitation \texttt{PRESS A} is added, at $x=132$,
which centres 7 characters of 8 pixels; in \texttt{PLAY} and \texttt{OVER}
the score, the worm and the food. The output depends only on the state and on
the memory: \code{draw} changes nothing. It is a Moore machine.

\asmfile[caption={Step 2 \textbf{(changed)}: the end of \code{verm1\_draw}, and \code{draw\_text}.},label=lst:v3:draw2]{code/l17/s19.S}

\texttt{OVER} adds two words on top of the scene
(Listing~\ref{lst:v3:draw2}). The helper \code{draw\_text} writes white on
black: it fills \reg{a3} and \reg{a4} and lets \reg{a0} to \reg{a2} pass
through, removing two lines from every call that writes text.

\subsection{Four functions that do nothing}

\asmfile[caption={Step 2: temporary stubs.},label=lst:v3:stubs]{code/l17/s20.S}

The stubs of Listing~\ref{lst:v3:stubs} make the names exist, so that the
program builds; \code{new\_game} only changes the state. Steps 3 and 4
replace them one at a time. It is the method of
Chapter~\ref{ch:firstcontact}: start with returns that do nothing, then add
one thing at a time.

\begin{prgcheck}[Run it: a machine with nothing inside]
\begin{enumerate}[nosep]
\item Build and run: the title line, and \texttt{PRESS A} in the middle of the
  screen.
\item Press \key{A}: the invitation disappears. A red cell is in the top-left
  corner: the food, at column 0 of row 0, because \code{food\_x} and
  \code{food\_y} start at zero and the stub of \code{new\_game} does not call
  \code{place\_food}.
\item Nothing moves and the game cannot end: \texttt{PLAY} has no rules yet.
\item In the debugger, find the variable \code{state} and watch it change from
  0 to 1 when you press \key{A}.
\end{enumerate}
\end{prgcheck}

% ----------------------------------------------------------------------------
\section{Step 3: a worm in memory}

\subsection{\texttt{new\_game}: empty the map, lay the worm}

Replace the stub of \code{new\_game} with Listings \ref{lst:v3:new1}
and~\ref{lst:v3:new2}.

\asmfile[caption={Step 3 \textbf{(new)}: \code{new\_game}, first half.},label=lst:v3:new1]{code/l17/s23.S}

The 960 bytes of the map are set to zero, one \ins{sb} at a time. Then three
segments are laid in columns 3, 4 and 5 of row 12, and each is written twice:
in the ring and in the map. The expression \code{grid+12*COLS} is computed by
the assembler: it is the address of the first byte of row 12. \ins{sb}
stores one byte; the index is added to the base address with no shift,
because each entry is one byte.

\asmfile[caption={Step 3 \textbf{(new)}: \code{new\_game}, second half.},label=lst:v3:new2]{code/l17/s24.S}

$\code{head\_i}=2$ and $\code{length}=3$ describe the worm; direction, timer
and score start from zero. The state is stored last: the transition from
\texttt{TITLE} to \texttt{PLAY} is complete only when everything is ready.

\subsection{\texttt{draw\_worm}: walk the ring backwards}

\asmfile[caption={Step 3 \textbf{(new)}: \code{draw\_worm}.},label=lst:v3:worm]{code/l17/s25.S}

Segment $i$ is the entry $(\code{head\_i}-i)$ AND 255
(Listing~\ref{lst:v3:worm}). \ins{lbu} loads one byte as an unsigned number,
and not \ins{lb}: a column is a number from 0 to 39, never negative. Both
would work here, since with \ins{lb} only a byte of 128 or more would become
negative (Chapter~\ref{ch:integers}); but \ins{lbu} says what is meant.
Segment 0 is the head, in a brighter green.

\subsection{\texttt{draw\_score}}

\asmfile[caption={Step 3 \textbf{(new)}: \code{draw\_score}.},label=lst:v3:score]{code/l17/s26.S}

This is the loop that was inside \code{verm1\_draw} in Part~2, now a function
(Listing~\ref{lst:v3:score}). It uses \reg{s0}, so it saves it and restores
it. One stub is left (Listing~\ref{lst:v3:stub2}): the worm exists, and
stands still.

\asmfile[caption={Step 3: the last stub.},label=lst:v3:stub2]{code/l17/s27.S}

\begin{prgcheck}[Run it: a worm that stands still]
\begin{enumerate}[nosep]
\item Build and run, press \key{A}: a worm of three cells in row 12, columns 3
  to 5, the head, brighter, on the right.
\item The food is now at a random cell: \code{new\_game} calls
  \code{place\_food}.
\item In the debugger, break at the start of \code{draw\_worm} and step until
  \code{la t0, body\_x} has run: \reg{t0} holds the address of
  \code{body\_x}. Type it in the \emph{Memory} field: the first bytes are
  \texttt{03 04 05}.
\item Find $\code{grid}+12\times 40+3$: three bytes at \texttt{01} in a sea of
  zeros. (\code{grid} is 512 bytes after \code{body\_x}, and row 12 starts 480
  bytes after \code{grid}.)
\end{enumerate}
\end{prgcheck}

% ----------------------------------------------------------------------------
\section{Step 4: rules}

One step of the worm has three movements: decide, look, write.

\subsection{\texttt{play\_update}: the update of Part 2}

Replace the last stub with Listings \ref{lst:v3:play1}
and~\ref{lst:v3:play2}. They are the body of the old \code{verm1\_update}
under a new name; the buttons arrive in \reg{a0}, so that there is no call to
read them.

\asmfile[caption={Step 4 \textbf{(changed)}: \code{play\_update}, the steering.},label=lst:v3:play1]{code/l17/s30.S}

\asmfile[caption={Step 4 \textbf{(changed)}: \code{play\_update}, the timer.},label=lst:v3:play2]{code/l17/s31.S}

\subsection{\texttt{step}: three registers to keep}

\asmfile[caption={Step 4 \textbf{(changed)}: \code{step}, the beginning.},label=lst:v3:s1]{code/l17/s32.S}

\code{step} needs three values that must survive a call to
\code{place\_food}: the cell of the new head, in \reg{s0} and \reg{s1}, and a
flag in \reg{s2} that is zero when the worm eats. They are kept in \reg{s}
registers, and the frame has room for four words: \reg{ra} and the three
saved registers (Listing~\ref{lst:v3:s1}). The table lookup is that of
Part~2.

\subsection{\texttt{step}: where will the head be?}

\asmfile[caption={Step 4 \textbf{(changed)}: \code{step}, the new cell.},label=lst:v3:s2]{code/l17/s33.S}

The present head is the entry \code{head\_i} of the ring: one byte each for
column and row, loaded with \ins{lbu}. \code{dx} and \code{dy} are added and
the result is wrapped as in Part~2 (Listing~\ref{lst:v3:s2}). Nothing is
written yet: first the program must look.

\subsection{\texttt{step}: eat, or move the tail}

\asmfile[caption={Step 4 \textbf{(changed)}: \code{step}, the tail.},label=lst:v3:s3]{code/l17/s34.S}

The comparator of Part~2 leaves its result in \reg{s2}
(Listing~\ref{lst:v3:s3}). If the worm eats, the length grows by one and the
tail stays where it is. If it does not, the tail is found at the entry
$\code{head\_i}-\code{length}+1$, and its column and row are put in \reg{t3}
and \reg{t4}. At 255 cells the worm stops growing.

\begin{prgconcept}
Growing and not moving the tail are the same thing. It is the whole trick of
Snake.
\end{prgconcept}

\subsection{\texttt{step}: free the tail, test the head}

\asmfile[caption={Step 4 \textbf{(new)}: \code{step}, the collision.},label=lst:v3:s4]{code/l17/s35.S}

The number of a byte of the map is $\mathit{row}\times 40+\mathit{column}$:
one \ins{mul} and one \ins{add} (Listing~\ref{lst:v3:s4}); \ins{mul} belongs
to the M extension of \riscv, which the ESP32-C6 has. The byte of the tail
is cleared \emph{first}, and then the byte of the head is tested. The order
matters: in this way the head may enter the cell that the tail has just left,
and the worm can chase its own tail, as in the original game. If the cell is
occupied the state becomes \texttt{OVER}, and \code{best} is updated.

\subsection{\texttt{step}: write the new head}

\asmfile[caption={Step 4 \textbf{(new)}: the end of \code{step}.},label=lst:v3:s5]{code/l17/s36.S}

The cell is marked in the map, while \reg{t0} still points to it; the index
\code{head\_i} becomes $(\code{head\_i}+1)$ AND 255, and the ring closes;
column and row are written in the new entry (Listing~\ref{lst:v3:s5}). If
\reg{s2} is zero there is a point, new food and a sound. After a collision
the function jumps straight to its last label: nothing is written, and the
screen shows the worm as it was before the fatal step.

\subsection{\texttt{place\_food}: not on the worm}

\asmfile[caption={Step 4 \textbf{(changed)}: \code{place\_food}.},label=lst:v3:food]{code/l17/s37.S}

A cell is drawn at random, and then the map is consulted: if the cell is
occupied another is drawn, until a free one comes
(Listing~\ref{lst:v3:food}). The column waits in \reg{s0} during the second
call. The loop ends quickly: even with 255 cells out of 960 occupied, the
average is fewer than two attempts.

\begin{prgcheck}[Run it: Verm1, Part 3]
The worm follows the head and grows (Figure~\ref{fig:v3:play}); turning into
the body ends the game (Figure~\ref{fig:v3:over}); \key{A} returns to the
title, and \key{A} again starts a new game.
\end{prgcheck}

% ----------------------------------------------------------------------------
\section{Step 5: look at memory move}

\begin{prgtry}[Watch the ring turn]
\begin{enumerate}[nosep]
\item Run in debug mode, press \key{A} to start a game, then pause.
\item Press \key{Update} and put a breakpoint on the first instruction of
  \code{step}.
\item \key{Resume}: the program stops once per step of the worm, not once per
  frame.
\item Type the address of \code{body\_x} in the \emph{Memory} field: the bytes
  \texttt{03 04 05} of the start, then \texttt{06}, \texttt{07}, \ldots, one
  new byte per step.
\item Go down for a few steps: \code{body\_x} repeats the same column, and
  \code{body\_y}, 256 bytes further on, counts the rows.
\item Eat, and compare \code{head\_i} and \code{length} before and after: both
  grow by one.
\end{enumerate}
\end{prgtry}

The bytes of the old positions are never erased: the ring keeps the whole
history of the last 256 steps, and \code{length} says how much of it is still
worm. After 256 steps the writing starts again from the first byte.

\subsection{Catch the edge of button A}

Put a breakpoint on \code{and s1, s0, t1} in \code{verm1\_update}, hold
\key{A} and press \key{Resume} three times. Table~\ref{tab:v3:edge} shows
what is read: \reg{s1} is \hex{10} in one frame only, the frame in which the
button goes down. (The test of the reference file holds \key{A} for 28 frames
and checks that exactly one transition happens.)

\begin{table}[htb]
\centering\small
\caption{The edge of button \key{A}.}\label{tab:v3:edge}
\begin{tabular}{@{}cccccc@{}}
\toprule
\textbf{Frame} & \textbf{A is} & \textbf{\reg{s0} (now)} & \textbf{\code{prev\_btn}} & \textbf{\reg{s1} = \reg{s0} AND NOT prev} & \textbf{Effect}\\
\midrule
$n$ & up & \hex{00} & \hex{00} & \hex{00} & --\\
$n+1$ & down & \hex{10} & \hex{00} & \hex{10} & transition\\
$n+2$ & down & \hex{10} & \hex{10} & \hex{00} & --\\
$n+3$ & up & \hex{00} & \hex{10} & \hex{00} & --\\
\bottomrule
\end{tabular}
\end{table}

\subsection{Three questions for the debugger}

\begin{itemize}
\item \emph{Where is \code{state}?} After \code{la t0, state} in
  \code{verm1\_update}, \reg{t0} holds its address. Watch \texttt{00},
  \texttt{01}, \texttt{02} as you play and lose.
\item \emph{Is the map consistent with the ring?} Pause, read the head from
  \code{body\_x} and \code{body\_y}, compute
  $\mathit{row}\times 40+\mathit{column}$, and look at that byte of
  \code{grid}.
\item \emph{What did the fatal step see?} Break on the \ins{lbu} that follows
  the second \ins{mul} in \code{step}: \reg{t1} is 1 exactly once, at the
  bite. At that moment \reg{s0} and \reg{s1} hold the cell that the head
  wanted to enter, and the screen shows that it is part of the body.
\end{itemize}

\begin{prgconcept}
A program is right when its memory is right. Learn to read it.
\end{prgconcept}

% ----------------------------------------------------------------------------
\section{Check your understanding}

\begin{quiz}{$\code{head\_i}=1$ and $\code{length}=4$. Which entries of the ring hold the worm?}
\item 1, 2, 3, 4
\item 0, 1, 2, 3
\item 254, 255, 0, 1
\item 253, 254, 255, 0
\end{quiz}

\begin{quiz}{Which byte of \code{grid} describes the cell in column 7 of row 3?}
\item byte 37
\item byte 127
\item byte 283
\item byte 307
\end{quiz}

\begin{quiz}{The player holds \key{A} down for two seconds on the title screen. What happens?}
\item the game starts and ends at once
\item the game starts once
\item nothing, until \key{A} is released
\item the states cycle at every frame
\end{quiz}

\begin{quiz}{Why is the tail's byte cleared before the head's byte is tested?}
\item to save an instruction
\item so that the head can enter the cell the tail leaves
\item because \ins{sb} is faster than \ins{lbu}
\item the order does not matter
\end{quiz}

\section{Make it yours}

\begin{exercise}[on \prg]
\begin{enumerate}[nosep]
\item Walls: the worm dies when it leaves the screen instead of wrapping.
\item Start with a worm of five cells; then let the length at the start depend
  on the best score.
\item Poison: a second, blue cell that makes the worm one cell shorter.
\item On the game-over screen show the best score as marks on the second line.
\end{enumerate}
\expected{Each change touches one or two functions. Keep a copy of the
working Part~3: Part~4 starts from it. (Hints: for the walls, replace each
wrap in \code{step} with a jump to the code that stores \texttt{OVER}. For
the poison, move the tail twice in one step, clearing both bytes of the map,
and never go below length 1.)}
\end{exercise}

\section{Two functions that did not change}

\code{eat\_sound} and \code{draw\_cell} are those of Part~2 and are kept as
they are. They are reprinted here so that this chapter contains the whole
program: the complete file is the sum of the listings marked new and changed,
plus these two.

\asmfile[caption={\textbf{Unchanged}: \code{eat\_sound}.},label=lst:v3:sound]{code/l17/s49.S}

\asmfile[caption={\textbf{Unchanged}: \code{draw\_cell}.},label=lst:v3:cell]{code/l17/s50.S}

\section{Where we are, and what comes next}

\vermone is now a game: it has a goal, a risk, a beginning and an end. A
machine with three states, a ring of bytes, a map: sequential logic and memory
written as a program, in about 450 lines of assembly that the reader has
written and can explain. Three things are missing: the program redraws the
whole screen at every frame, the score is a row of marks, and the speed never
changes. The next four chapters open the machine that runs those lines: how
the processor executes, calls and division, how video and sound work. Part~4
is the finished game.
